\chapter*{RESUMO}
\addcontentsline{toc}{chapter}{RESUMO}

A reprodutibilidade é um pilar do método científico, mas a pesquisa em robótica móvel
raramente a alcança de forma sistemática. Configurar um experimento de navegação,
registrar dados sensoriais ao longo de múltiplas execuções e comparar as trajetórias
resultantes com rigor quantitativo exige uma infraestrutura que a maioria dos laboratórios
constrói de forma improvisada, específica para cada projeto, sem possibilidade de reuso.
Esta dissertação apresenta um framework que torna essas etapas explícitas, automatizadas
e transferíveis.

O sistema proposto --- denominado \textit{Fleet UI} --- orquestra experimentos de navegação
autônoma com um ou mais robôs móveis operando sob ROS~2 Jazzy e a pilha de navegação Nav2.
Uma arquitetura baseada em papéis separa as unidades em \textit{MUUT} (Mobile Unit Under
Test), responsável pela execução e gravação das rotas, e \textit{FUUT} (Fixed Unit Under
Test), que observa o cenário de uma posição fixa. Essa separação permite que o mesmo
testbed físico sirva a diferentes desenhos experimentais sem modificação da camada de software.

O protocolo de gravação--reprodução constitui o núcleo experimental: uma rota é definida
por navegação sequencial de waypoints; o orquestrador armazena as poses atingidas; as
execuções subsequentes de replay reproduzem essas poses via a ação \texttt{FollowWaypoints}
do Nav2. Bags de sensores no formato MCAP capturam LiDAR, odometria de roda, IMU e
pose estimada pelo SLAM em cada execução. A análise \textit{offline} calcula o
\textit{Root Mean Square Error} (RMSE) entre trajetórias reamostradas em tempo
normalizado, gerando sobreposições visuais e métricas de repetibilidade.

A validação quantitativa, realizada em simulação Gazebo Harmonic com uma única
unidade TurtleBot4 Standard, utilizou uma baseline limpa e dez execuções independentes
de replay. Cada replay foi iniciado a partir do mesmo estado por relançamento completo
da simulação, da pilha Nav2 e dos nós de coleta, garantindo que a odometria registrada
começasse praticamente no mesmo ponto em todas as execuções. A campanha obteve RMSE
médio de 3,35\,cm entre cada replay e a execução de referência, com máximo de
4,76\,cm, em trajetórias com comprimento médio de 0,826\,m. O erro final médio foi de
1,64\,cm (máximo de 3,26\,cm), e a duração média foi de 17,51\,s com desvio-padrão de
0,36\,s. Esses resultados ficam confortavelmente abaixo da tolerância padrão do Nav2 de
25\,cm; a arquitetura suporta múltiplas unidades, mas a avaliação empírica de cenários
com frota simultânea não foi conduzida nesta dissertação. Uma interface web em React e
FastAPI permite que pesquisadores definam experimentos, monitorem o estado dos sensores
e visualizem trajetórias sobre o mapa SLAM sem exposição direta à linha de comando do
ROS~2.

\vspace{0.5cm}
\noindent\textbf{Palavras-chave}: navegação autônoma; repetibilidade de trajetórias;
frota de robôs móveis; ROS~2; Nav2; SLAM; coleta de dados sensoriais; framework experimental.
